%% ============================================================
%% MODELO DE ATIVIDADE · GENÉRICO
%%
%% Esta é a primeira atividade. O melhor é trocá-la pelo modelo do
%% tipo certo (modelos/pcb.tex, modelos/simulacao.tex…): ele já traz a
%% ficha técnica e as perguntas do tipo. Os tipos estão em nro/tipos.tex. Cada \preencher{…} é uma pendência: troque-o pelo seu
%% texto. Cada \figuraprovisoria também: troque-a pela figura de verdade.
%%
%% Uma atividade é UMA entrega sua, com começo, meio e fim. "Participei
%% do projeto X por dois anos" não é uma atividade: são várias. Divida.
%% ============================================================
\begin{atividade}{O que foi feito, em poucas palavras}{
  tipo         = {},   % o tipo: pcb, mecanico, simulacao, software, ensaio, clinico, pesquisa,
                       % extensao, ensino, evento, gestao, comunicacao ou outro
  modalidade   = {},   % a chave da modalidade pedida: EXT, PES, COMP, VPC, EVT, PRD, ORG, ENS, GES, CUR, OUT
  alternativa  = {},   % a modalidade alternativa, se o Colegiado não aceitar a primeira
  horas        = {},   % as horas desta atividade, com base nos registros (Parte IV)
  inicio       = {},   % AAAA-MM
  fim          = {},   % AAAA-MM, ou: em andamento
  projeto      = {},
  papel        = {},   % Responsável, Corresponsável ou Colaborador(a)
  supervisor   = {},   % quem confirma: nome e cargo (o supervisor do projeto, o gerente)
  registros    = {},   % os códigos no SOMA, com ponto e vírgula: NBL-14; NRO-PRO-003-7
  competencias = {},   % as competências das DCN que ela desenvolveu (I a VIII): I, III, V
  disciplinas  = {},   % as disciplinas do seu curso ligadas a ela: ELE075 Eletrônica; ELE067 Circuitos
}

\begin{contexto}
\preencher{O que foi feito, com que objetivo, e por que isso é formação no seu curso.}
\end{contexto}

\begin{contribuicao}
\preencher{O que foi seu, separado do que foi da equipe.}
\end{contribuicao}

\begin{desenvolvimento}
\preencher{Como o trabalho foi feito, com quem, com que método e com que ferramentas.}
\end{desenvolvimento}

\begin{resultados}
\preencher{O resultado, com números quando houver, e como ele pode ser conferido.}
\end{resultados}

\begin{evidencias}
% Uma linha por evidência: \evidencia{tipo}{identificador}{o que mostra}{onde conferir}
% \evidencia{<tipo>}{<identificador>}{<o que mostra>}{<onde conferir>}
\end{evidencias}

\begin{aprendizados}
\preencher{O que você aprendeu, e o que do curso você aplicou.}
\end{aprendizados}
\end{atividade}
